\hspace{3mm}

\centerline{\bf \Huge Турниры}

\centerline{\bf \large Базовый уровень}

\begin{flushright}
        \scriptsize{}
\end{flushright}

\problem{1}
а) Сколько шахматистов играло в однокруговом турнире, если всего в этом соревновании было сыграно 190 партий?

\noindent
б) Сколько шахматистов играло в однокруговом турнире, если всего было набрано больше 50, но меньше 60 очков?

\noindent
в) Сколько команд играло в однокруговом турнире, если всего было набрано больше 50, но меньше 60 очков по системе 2–1–0?

\problem{2}
В шахматном турнире каждый участник выиграл белыми столько же партий, сколько все остальные вместе чёрными. Докажите, что у всех поровну побед.

\problem{3}
Шахматист сыграл в турнире 20 партий и набрал 12,5 очков. На сколько больше он выиграл партий, чем проиграл?

\problem{4}
В однокруговом турнире была применена система подсчёта очков 2–1–0. Вася набрал очков меньше Пети. Может ли у него стать очков больше, чем у Пети, если результаты пересчитать

а) по шахматной системе (за победу 1 очко, за ничью 1 , 2 за поражение 0),

б) по футбольной системе (за победу 3 очка, за ничью 1,
за поражение 0)?

\problem{5}
В однокруговом турнире десяти команд команда А набрала очков больше любой другой, а Я — меньше любой другой. Какой наименьший разрыв в очках может быть между командами А и Я, если очки считаются

а) по системе 2–1–0,

б) по футбольной системе?



\newpage

\hspace{3mm}

\centerline{\bf \Huge Турниры}

\centerline{\bf \large Продвинутый уровень}

\begin{flushright}
        \scriptsize{}
\end{flushright}

\problem{1}


\problem{2} 


\problem{3} 

\problem{4}

\problem{5} 
